\begin{frame}{학습 곡선: Block B에서 x축이 ``에포크''인 이유}
  \begin{itemize}
    \item 8.1: x축 = $n\_estimators$ — GBDT는 모델을
      \textit{크게} 만드는 파라미터가 트리 수이므로 x축이 곧
      \textbf{용량}
    \item Block B: 신경망은 \textbf{구조(층 수·너비)로 용량을
      고정}하고, \textbf{학습량(에포크)을 바꿔} train/val 곡선
      을 그리는 것이 표준
    \item 구조를 바꾸면 모델 \textbf{자신이 달라지므로}(트리 수
      만큼 ``같은 모델의 크기''가 아님) \textbf{에포크를 x축}
      으로
  \end{itemize}
  \vskip0.4em
  \small
  \begin{tabular}{@{}lcc@{}}
    \toprule
    GBDT ($n\_estimators$) & train & val \\
    \midrule
    20 & 0.988 & 0.944 \\
    50 & \textbf{1.000} & 0.961 \\
    200 & 1.000 & 0.983 \\
    \bottomrule
  \end{tabular}
  {\footnotesize train은 50트리에서 \textbf{1.0 포화}.}
\end{frame}
